\documentclass[12pt]{article}
\usepackage[margin=1in]{geometry}
\usepackage{amsmath,amssymb,amsthm,mathtools}
\usepackage{booktabs,threeparttable,graphicx,setspace,enumitem,longtable,array}
\usepackage{xcolor}
\usepackage{etoolbox}
\AtBeginEnvironment{tabular}{\footnotesize}
\usepackage[authoryear,round]{natbib}
\usepackage[colorlinks=true,linkcolor=blue!50!black,citecolor=blue!50!black,urlcolor=blue!50!black]{hyperref}
\usepackage{fancyhdr}
\doublespacing
\emergencystretch=2em
\pagestyle{fancy}\fancyhf{}
\fancyhead[L]{\small Niroula, Niroula and Khatiwada}
\fancyhead[R]{\small What Prices Reveal About Hidden Traders}
\fancyfoot[C]{\thepage}
\renewcommand{\headrulewidth}{0.3pt}
\setlength{\headheight}{14pt}

\newtheorem{theorem}{Theorem}[section]
\newtheorem{proposition}[theorem]{Proposition}
\newtheorem{lemma}[theorem]{Lemma}
\newtheorem{corollary}[theorem]{Corollary}
\theoremstyle{definition}
\newtheorem{definition}[theorem]{Definition}
\newtheorem{assumption}[theorem]{Assumption}
\theoremstyle{remark}
\newtheorem{remark}[theorem]{Remark}
\newcommand{\E}{\mathbb{E}}
\newcommand{\R}{\mathbb{R}}
\newcommand{\KL}{D_{\mathrm{KL}}}
\newcommand{\lab}[1]{\textsc{[#1]}}

\title{What Prices Reveal About Hidden Trader Populations,\\ and What Identifying Them Costs}
\author{Pratik Niroula\thanks{Department of Computer Information Science, Minnesota State University, Mankato, MN 56001, USA. \texttt{pratik.niroula@mnsu.edu}.}
\and Prashant Niroula\thanks{Faculty of Management, Tribhuvan University, Kirtipur, Nepal. \texttt{prashantniroula20@gmail.com}.}
\and Sujan Khatiwada\thanks{Department of Electrical, Computer and Robotics Engineering, Minnesota State University, Mankato, MN 56001, USA. \texttt{sujan.khatiwada@mnsu.edu}.}}
\date{Development revision, 30 October 2026}

\usepackage{array}
\begin{document}
\maketitle
\thispagestyle{empty}
\begingroup\onehalfspacing

\begin{abstract}
\noindent Can market prices tell apart groups of traders who react to past prices at different speeds? In a linear market in which groups follow prices with exponentially fading memory and their combined orders move the price, prices reveal only one combined summary of the groups. Groups with different memories can be told apart, but the cost rises steeply as their memories approach: for $K$ unknown groups whose memories are $\Delta\eta$ apart, the number of controlled experiments grows as $(\Delta\eta)^{-(4K-2)}$, with matching bounds on suitably rich classes, and watching without trading needs at least as many observations. For a known pair, the trading capital needed rises sixteen-fold each time their difference is halved. Records that report each group's flow separately keep the problem well conditioned, and a stylised supervisory decision shows when that matters. In pre-specified tests, published trader-class positions in eight CFTC futures markets, and Binance class data in 29 of 29 new markets, pin down a position persistence that prices leave open, and the predicted rates hold with real Bitcoin price noise where it can be used directly. We do not show that this persistence is learning memory, and several pre-specified tests failed, including a test of what prices do reveal and a test of CFTC class flows.
\end{abstract}

% BEGIN ADVANCED FRONT STATUS 20261030
\paragraph{Research status update: 30 October 2026.}
This development revision records a verified returned synthetic experiment and
a dated work register at the end. The new audit supports a conditional
finite-horizon financial decision result, not general profitability or
real-market safety. Earlier theoretical and historical claims are retained
with their stated assumptions; this update does not independently verify every
proof or rerun every earlier empirical study. Novelty, full proof review,
economic information acquisition and the public-data bridge remain open.
% END ADVANCED FRONT STATUS 20261030


\noindent\textbf{Keywords:} algorithmic trading, identification, trader heterogeneity, market microstructure, system identification, minimax sample complexity, value of disclosure, audit trails.\\
\textbf{JEL:} C14, C32, C58, D83, G14, G18.

\paragraph{Reproducibility.} Every number in the paper is produced by a script or notebook cell that writes a result file from public data or from a closed form; Appendix~\ref{app:repro} maps each table, figure and number to its source. ``Registered'' or ``pre-specified'' in this paper means that the authors wrote the specification and committed it to the public repository before opening the test data; this is the authors' declaration, not independent preregistration. For the October 2026 follow-up tests (Section~\ref{sec:learners} and Section~\ref{sec:disc_data}) each specification reached the public repository in the same push as its results, so the order rests on the commit history alone. For the real-noise test, the class-flow test and the CFTC robustness checks (Sections~\ref{sec:injection} and~\ref{sec:cot}), specifications and code were pushed in a separate commit (\texttt{84af69f}, 03:43 UTC on 31 October 2026) before the first of them ran (03:56 UTC). The proofs have been checked by the authors and against numerical computations, and they have not yet been independently reviewed; the numerical checks reproduce results with the same code unless stated otherwise. Lean~4 with Mathlib machine-checks, using only Lean's standard axioms, two of the main results: the equivalence in Theorem~\ref{thm:merged}(i) for the model itself (for known $a$ and $\kappa\neq0$, two economies have the same price response if and only if they have the same merged measure), and the lower bound of Theorem~\ref{thm:kcohort} on a finite observation horizon, from the moment gap through the closed loop and the Gaussian divergence to the testing step. For two groups the check is complete; for general $K$ it assumes that a moment-matched pair of $K$-point rules exists (the Gauss and Gauss--Radau rules supply one). The upper bound (Theorem~\ref{prop:esprit}), the general-memory and passive-data theorems and the central limit step of Proposition~\ref{prop:tagged} are not formally verified. The results whose proofs most need an independent check are Theorems~\ref{prop:esprit}, \ref{thm:general}, \ref{prop:passive}, \ref{thm:spectral} and~\ref{thm:rh}. Numerical constants such as the capital bounds are computed at a declared parameter set and are not separate theorems.

\endgroup
\newpage
\setcounter{page}{1}
\tableofcontents
\newpage

\section{Introduction}

Markets are increasingly run by algorithms. Trend-following funds, market-making programs, execution algorithms and, more recently, AI trading agents all read recent prices, decide what to trade, and by trading move the prices they read. When something goes wrong, a sharp sell-off, a sudden loss of liquidity, a crowded trade unwinding, the first questions supervisors and risk managers ask are who was trading and how. Which group of traders drove the move? How quickly do they react, and how long do they remember what prices did? Could the same group do it again?

This paper asks how much of that can be learned from what everyone can see, prices and anonymous order flow, and, when public data are not enough, what it costs to find out in other ways.

The question matters now. Trading is passing to learning algorithms and AI agents \citep{colliard2022,dou2025}, and regulators have begun to require firms to say which algorithm sent each order. Under MiFID II, EU trading venues must keep order records that identify the person or algorithm responsible for each investment and execution decision \citep{rts24,esma2016}. The US Consolidated Audit Trail (CAT) is designed to record each order in US-listed stocks and options, and the account behind it, across exchanges and broker-dealers; its revised 2026 budget, published in March 2026, is about \$147 million \citep{sec2012,catbudget2026}. Such systems serve many purposes, among them enforcement, rebuilding market events after the fact, and surveillance. One question they raise, and the one we study, is how much of the information about who drives price moves could instead be learned from public data, and at what cost.

\paragraph{The short answer.} Prices tell an observer less than one might hope, and sharpening what they do tell is expensive.
\begin{itemize}[itemsep=1pt]
\item Prices reveal one combined summary of all the trader groups, not the groups themselves. Two groups that react to prices with the same memory look exactly like one larger group, however much price data is collected (Theorem~\ref{thm:merged}).
\item Groups with different memories can be told apart in principle, but the data needed grow very fast as their memories get closer. For two unknown groups, halving the difference between their memories multiplies the data needed by 64; for three groups, by 1{,}024 (Theorems~\ref{thm:kcohort} and~\ref{prop:esprit}).
\item Records that report each class's flow separately remove this problem. Each group then becomes an ordinary estimation problem that stays well conditioned however similar the groups are (Proposition~\ref{prop:tagged}).
\end{itemize}
We keep three statements apart throughout: what cannot be learned at all, what can be learned but only at high cost, and where estimates become unstable.

\paragraph{Why prices hide the groups.} The reason is the feedback loop. Algorithms trade on recent prices, their orders move prices, and those prices feed back into the algorithms. We model each group as tracking the price with a fading memory: recent prices count most, older prices count less, and one number, the group's memory rate, sets how fast the past fades. A desk behaves exactly this way if it uses a Kalman filter or a regularized online learning rule (Lemma~\ref{lem:micro}). Because every group reaches the price only through the sum of all their trades, the price records the total and loses the parts. Think of several people pushing the same door: from the door's movement one can learn the total push and how it builds and fades, but two people who push in the same rhythm look like one stronger person.

\paragraph{Three ways to learn more, and what each costs.} An observer that wants to know more, whether a supervisor, an exchange's surveillance team or a large trading desk, has three options (Section~\ref{sec:observer}).
\begin{enumerate}[label=(\alph*),itemsep=1pt]
\item \emph{Wait and watch prices.} This costs time. In a stationary Gaussian market observed without error, a passive observer needs at least as many observations, in order of magnitude, as an active one needs experiments (Theorem~\ref{prop:passive}).
\item \emph{Trade against the groups} and watch how the price responds, much as one taps a wall to find a hidden beam. Every test trade moves the price against the trader, so learning this way costs money. To tell one group from two groups whose memory rates differ by $2\delta$, the quantity traded grows like $1/\delta^{2}$ and the capital at risk like $1/\delta^{4}$: halve the difference and the capital needed rises sixteen-fold (Theorems~\ref{thm:quantity} and~\ref{thm:cap}). If memory is a general mixture of fading memories rather than a few exponentials (with the slowest bounded away from a random walk), accuracy improves only with the logarithm of the data, so halving the error squares the data needed (Theorem~\ref{thm:general}).
\item \emph{Acquire labelled records}, such as audit-trail data or class-level position reports. This costs whatever the reporting system costs, but it removes the explosion in data: labels save a factor of order $(\Delta\eta)^{-(4K-2)}$ at fixed accuracy, where $K$ is the number of groups and $\Delta\eta$ the spacing of their memory rates (Theorem~\ref{thm:discvalue}).
\end{enumerate}
Proposition~\ref{prop:choice} says when each option is the cheapest.

\paragraph{How large are the costs? A hypothetical illustration.} To give a sense of scale we price the trading option in a hypothetical benchmark. The model parameters are illustrative; only the size of the measurement noise is estimated, from Bitcoin data. An outside participant that re-learns two unknown algorithms each week with test trades would lose at least \$15 million to \$356 million a year per market; with three unknown algorithms it would lose more than \$40 million a year per market even when their memories are far apart (Section~\ref{sec:disc_dollars}). These are lower bounds for one stated way of learning, not estimates of what any regulator spends: regulators do not trade to learn. They are best read as the price of the information that labelled records deliver without any trading.

\paragraph{Why this matters for finance.} The results point in a constructive direction for several groups of readers.
\begin{itemize}[itemsep=1pt]
\item \emph{Supervisors and policymakers} get a way to decide what data to require. If the goal is only to watch overall price risk, public prices are enough, because the combined price response is learned at the ordinary rate (Remark~\ref{rem:riskneutral}). Labelled records earn their cost when supervision must be targeted, for example when one algorithm must be paused while the others keep supplying liquidity, and when the reporting system covers enough markets to share its fixed cost (Proposition~\ref{prop:choice}). In our stylised cost model this illustrates a possible economy of scale for market-wide reporting; it does not explain the legal design of regimes such as CAT. For the narrow task of attributing behaviour to classes, a persistent pseudonymous label is enough, which bears on how much personal data such a task needs but not on what enforcement needs (Section~\ref{sec:policy}). Proposition~\ref{prop:target} gives a stylised decision in which this matters: the size of a market-wide measure depends only on the combined response, while deciding whether halting one group is enough needs that group's own intensity and inherits the full cost of separating the groups.
\item \emph{Exchanges and trading venues}, which already see account-level flow, learn which questions about their participants public prices cannot answer, and why publishing class-level positioning, as some crypto venues already do, adds information that prices do not contain.
\item \emph{Risk managers and trading desks} get a realistic limit on what can be inferred about other traders or about crowding from market data, and a warning that common feedback measures computed from public data mostly reflect filtering and timing rather than true feedback (Section~\ref{sec:empirics}).
\item \emph{Researchers} who estimate models with several trader types get one reason why the parameters that separate the types are often imprecise, and a number that says how precise they can be before any estimation (Section~\ref{sec:lit}).
\end{itemize}

\paragraph{What the data show.} Two real reporting systems let us check the main prediction, that class records pin down what prices leave open. The U.S.\ Commodity Futures Trading Commission collects positions from large futures traders and publishes a weekly snapshot of them, aggregated by trader category. In a pre-specified test in eight futures markets (CME Bitcoin and Ether, the E-mini S\&P 500, 10-year Treasury notes, the euro, the yen, crude oil and natural gas), these class positions pin down the persistence of the designated speculative class's positions to a range 9 to 71 times narrower than weekly prices allow, while prices alone leave it almost anywhere in the admissible range (Section~\ref{sec:cot}). On Binance, in a pre-specified test on 30 markets never used before that matches the memory horizon of the two routes, the class data give the narrower interval in 29 of 29 evaluable markets, by a median factor of 7.3 (Section~\ref{sec:disc_data}).

\paragraph{Robustness, a new failure and real noise.} Rerunning the CFTC test with CME futures prices, without contract-roll weeks or with longer bootstrap blocks leaves its primary result unchanged. Read as weekly flows rather than positions, the same reports failed a second pre-specified test: they pin down how each class's flow responds to prices, but not how persistent that flow is. A third pre-specified test gives synthetic groups real Bitcoin price noise instead of Gaussian noise and finds the predicted rates, a slope of $-6.01$ against $-6$ for two groups and $-9.81$ against $-10$ for three, while class-tagged requirements stay flat; at large sample sizes it uses a normal approximation to the real noise (Section~\ref{sec:injection}). These comparisons are descriptive: a position series reveals its own persistence directly, while the price route must infer it. Several predictions failed, and Table~\ref{tab:scorecard} lists them next to the ones that passed. The data therefore show what prices cannot pin down. They do not show that this persistence is the learning memory of the model.

\begin{table}[!htbp]\centering\footnotesize
\begin{threeparttable}
\caption{What the empirical tests found: passes and failures}\label{tab:scorecard}
\setlength{\tabcolsep}{3pt}
\begin{tabular}{>{\raggedright\arraybackslash}p{0.27\textwidth}>{\raggedright\arraybackslash}p{0.38\textwidth}>{\raggedright\arraybackslash}p{0.27\textwidth}}
\toprule
Test & Prediction & Result\\
\midrule
CFTC, 8 markets, H1 & Class positions pin down a persistence that weekly prices leave open & Pass, 8 of 8\\
CFTC, H2 & The designated speculative class trend-follows & Fail, 4 of 8 (three markets have the opposite sign)\\
CFTC, H3 & That class has the less persistent positions & Pass, 6 of 8 (two indistinguishable)\\
CFTC flows (weekly changes), H1 & Class flows pin down their own persistence (95\% interval narrower than 0.5) & Fail, 3 of 8 (6 required); flow response nonzero in 6 of 8\\
Real-noise injection, $K=2$ & Price-only requirement grows as $(\Delta\eta)^{-6}$; tagged requirement stays bounded & Pass: slope $-6.01$ (large-$N$ points use a normal approximation); tagged max/min 2.8 ($\le8$ required)\\
Same test, $K=3$ (secondary) & Price-only requirement grows as $(\Delta\eta)^{-10}$ & Holds: slope $-9.81$\\
Binance, 30 fresh markets, (i)--(ii) & Class data restrict persistence; price-only intervals wide & Pass, 30 of 30 (one-hour price kernel)\\
Binance, same test, (iii) & Prices pin down the combined response & Fail, 5 of 30; the test failed its 80\% criterion\\
Binance, 30 new markets & Horizon-matched: class interval narrower & Pass, 29 of 29 evaluable (median factor 7.3)\\
Binance, same test, secondary & Class persistence is a stable memory & Not supported: close to a random walk after bias correction\\
Binance, 105-day follow-up & Combined response far cheaper to learn than the separate rates & Fail, 10 of 30 (23 required)\\
Binance, 58-day follow-up & Tagging advantage grows with the sample; parametric shrinkage & Pass, 22 of 29; fail, 20 of 29 (22 required)\\
Binance, return forecasts & Class memory states improve 30-minute return forecasts & Fail ($\Delta R^2=-1.09$ points, inconclusive)\\
Binance, state tracking & Class data track class positions better than prices & Pass, but explained by persistence of the target\\
\bottomrule
\end{tabular}
\end{threeparttable}
\end{table}

\paragraph{Two supplementary results.} Neither is needed for the main results. Markets with the same eigenvalues, the same one-step volatility and the same long-run price variance can still have very different worst-case short-run losses, so watching eigenvalues or volatility cannot replace knowing who trades (Theorem~\ref{thm:blind}). And standard feedback measures computed from public data mostly reflect filtering and clock alignment: two fee changes that moved trading volume by factors of 2.6 and 4.9 left measured feedback unchanged (Section~\ref{sec:empirics}).

\paragraph{Contribution.} A large literature asks how transparency affects liquidity and trading costs \citep{pagano1996,boehmer2005,bessembinder2006}. We ask a different question, what disclosure lets an observer learn about the traders behind price moves, and make three contributions. First, we show that a closed-loop price process identifies only the combined distribution of the trader groups' memories and weights, and nothing finer (Theorem~\ref{thm:merged}). Second, we price what it takes to go further under stated observation contracts, in experiments, observations, executed quantity and trading capital, with matching upper and lower bounds where we can prove them (Sections~\ref{sec:resources} and~\ref{sec:mixtures}). Third, we measure the value of class-level records in the same units and test, with published class positions, the part of the prediction that public data can test, reporting the failures alongside the passes (Section~\ref{sec:disclosure}). The mathematical tools, moment matching, two-point testing and Prony-type estimation, are known from mixture estimation and signal processing; what is new is their closed-loop, passive and financial form (Section~\ref{sec:lit}).

\paragraph{How to read the paper.} Every main result is stated formally and then explained in a short paragraph headed ``In words'', which says what it means for a trader or a supervisor; proofs are in the online supplement. Section~\ref{sec:worked} works through two desks that use Kalman filters, in numbers, and Table~\ref{tab:map} maps each question to its answer. Readers who want the finance first can read this introduction, Section~\ref{sec:observer}, the worked example and Section~\ref{sec:disclosure}, and take the statistical results as given. Table~\ref{tab:glossary} explains the main symbols.

\begin{table}[!htbp]\centering\small
\begin{threeparttable}
\caption{The main symbols in plain words}\label{tab:glossary}
\begin{tabular}{>{\raggedright\arraybackslash}p{0.17\textwidth}>{\raggedright\arraybackslash}p{0.75\textwidth}}
\toprule
Symbol & Meaning\\
\midrule
$K$ & Number of trader groups (``cohorts'')\\
$\eta_i$ & Group $i$'s learning rate: the weight it puts on the newest price. Large $\eta_i$ means a short memory and fast reactions; small $\eta_i$ means a long memory\\
$r_i=1-\eta_i$ & Group $i$'s memory rate: the share of last period's view that survives\\
$q_i$ & Group $i$'s trading intensity: how strongly it trades on its view. Positive $q_i$ means it buys when the price is above value, as trend followers do\\
$\Delta\eta$, $2\delta$ & How far apart the groups' learning rates are ($\Delta\eta$ for unknown groups, $2\delta$ for a known pair)\\
$a$, $\kappa$ & How persistent pricing errors are, and how far one unit of net order flow moves the price\\
$M$ (merged measure) & The combined summary of all groups that prices reveal: where the memories are and how much trading sits at each\\
$N$ & Number of experiments (test trades from a fresh start) or observations\\
Par units & The model's own unit of quantity and money, before conversion to USDT (Table~\ref{tab:capital})\\
\bottomrule
\end{tabular}
\end{threeparttable}
\end{table}

\paragraph{What we do not claim.} The identification results hold for the stated linear model and for the stated observation contracts, that is, for what the observer is assumed to see. The model treats price impact and the persistence of pricing errors as known primitives rather than solving a dealer equilibrium, and it assumes the observer knows the fundamental value; real data meet neither assumption. They are not general limits on surveillance with participant-level records, which can be far more informative (Propositions~\ref{prop:disclosure} and~\ref{prop:tagged}). The matching upper bound holds on a class with separated rescaled rates and weights bounded away from zero. Capital figures depend on the observation-noise calibration, which we estimate rather than assume. The empirical results are reduced-form diagnostics; they do not identify structural learning parameters. Our pre-specified real-data test of the identification boundary failed its criterion on one of its three parts: on 42 days of data the merged response was precisely estimated in only 5 of 30 markets (Section~\ref{sec:disc_data}). Binance's class ratios are positions, not tagged order flow, so they are the closest public analogue to the disclosure of Proposition~\ref{prop:tagged}, not an instance of it. We hold the algorithms' memories and intensities fixed across disclosure regimes; disclosure that reveals positions may lead informed desks to trade less or to split orders across accounts, which we do not model.

\subsection{Related literature}\label{sec:lit}

\paragraph{Algorithms, AI and supervision.} Recent work in finance asks how learning algorithms and AI agents behave in markets: reinforcement-learning market makers and informed traders can sustain wide spreads or collusive outcomes \citep{colliard2022,dou2025}, prices formed against Q-learners have distinctive dynamics \citep{banerjee2026}, and algorithmic markets can be manipulated \citep{tremacoldi2025}. Traders who differ in horizon or speed are studied in equilibrium by \citet{cespa2015} and \citet{biais2015}, and \citet{crouzet2020} model investors who trade at different frequencies. These papers model or simulate the algorithms. We take the supervisor's side and ask what can be learned about them from market data, and what disclosure is worth. A related literature measures the cost of regulation and of regulatory reporting \citep{trebbi2024}; our dollar comparison with the Consolidated Audit Trail is a statement about the benefit side: what the audit trail buys that prices cannot.

\paragraph{Price formation and impact.} Our price equation is a discrete linear impact model in the tradition of \citet{kyle1985}, \citet{glosten1985} and \citet{hasbrouck1991}, in which order flow moves prices and prices revert toward value. Transient impact and the cost of trading follow \citet{almgren2001}, \citet{obizhaeva2013}, \citet{gatheral2010} and \citet{bouchaud2004}; \citet{gatheral2010} shows that absence of dynamic arbitrage requires a positive-definite decay kernel, \citet{alfonsi2012} that excluding transaction-triggered manipulation requires convex, nonincreasing decay, and \citet{alfonsi2016} extend positive definiteness to several assets. Positive definiteness is the property Lemma~\ref{lem:toeplitz} delivers in our cohort economy. \citet{bouchaud2018} document order splitting and long memory in order flow, which is the main confounder in Section~\ref{sec:empirics}. Relative to this literature we treat the impact kernel not as a primitive but as the closed-loop output of adaptive traders, and ask what it reveals about them.

\paragraph{Algorithmic trading and feedback.} Algorithms that react to recent prices can amplify or dampen moves. \citet{brunnermeier2009} and \citet{shleifer1992} study liquidity spirals and fire sales, and \citet{kirilenko2017} document the behaviour of automated intermediaries in the 2010 Flash Crash. \citet{sornette2023} show that non-normal interactions can generate transient bubbles. Our contribution is on the measurement side: we show what a closed-loop observer can and cannot learn about the feedback, and that standard reduced-form indices are dominated by filtering and timing.

\paragraph{Learning algorithms and what can be inferred about them.} Empirical game-theoretic analysis treats the choice among simulated trading strategies as a game \citep{wellman2006}. Empirically, algorithmic trading has been found to improve liquidity and price discovery on average \citep{hendershott2011,brogaard2014}. These studies simulate the algorithms or observe them through regulatory data; we ask what an outside observer can infer about them from market data alone, and answer in samples and in capital. Our stance is that of partial identification \citep{manski2003}: when the structure is not identified, characterize what the data do determine, here the merged measure and, through it, the intensity-weighted distribution of the desks' information quality (Proposition~\ref{prop:primitives}).
\paragraph{Control without identification.} We do not claim that acting on a system without identifying it is new. \citet{moallemi2008} study strategic execution against an observer who learns from prices, which bears directly on our assumption that cohorts do not react to being observed, and \citet{parkvanroy2015} study execution that learns price impact while trading, with regret guarantees; our object is different, the latent trader cohorts behind the impact rather than the aggregate impact coefficients. Active identification of linear systems has matching bounds and optimal experiment design \citep{wagenmaker2020}, and task-directed exploration asks what must be learned for a decision \citep{wagenmaker2021}; our information matrix differs because the unknowns are zeros of a closed-loop response and the input is paid for in trading losses. \citet{vanwaarde2020} characterize when control is possible without unique system identification. Our results concern a specific market observation contract and its costs in executed quantity, capital and money.
\paragraph{Closest work.} The closest literature is super-resolution and moment-based mixture estimation, which gives rates for recovering clustered point sources and mixing measures from noisy moments \citep{moitra2015,heinrich2018,batenkov2021,liu2023,wu2020,kong2017}. Our $K$-cohort rate is, in that language, the square of the amplitude condition number $(\Delta\eta)^{-(2K-1)}$ for $K$ clustered nodes \citep{heinrich2018,wu2020}, and we do not claim the exponent as new. What changes when the mixture sits inside a closed-loop financial market is fourfold (Table~\ref{tab:novelty}): first, in mixture estimation and open-loop system identification the unknown rates are locations or poles of the observed law, whereas in a closed-loop market the cohorts feed back into the price, making the memory rates transmission \emph{zeros} of the observable transfer function $H(z)$ that must first be extracted from the loop (Theorem~\ref{thm:merged}); with $a$ and $\kappa$ known this extraction is a smooth invertible map, so it decides which method is valid (pole fitting on prices is not) rather than the rate; second, passive price observations face at least the same $(\Delta\eta)^{-(4K-2)}$ sample-complexity barrier as active probing, uniformly in the spacing, because $H(z)$ is minimum phase and the Whittle likelihood inherits the moment gap (Theorems~\ref{thm:minphase} and~\ref{prop:passive}); third, in a market statistical sample size translates into executed quantity $\Theta(\delta^{-2})$ and round-trip trading loss or margin capital $\Theta(\delta^{-4})$ via the self-impact Toeplitz quadratic form $\tfrac12 v^\top T_h v$ (Theorems~\ref{thm:quantity}--\ref{thm:rh}); and fourth, class-tagged disclosure breaks the feedback loop cohort by cohort and cuts the sample requirement by $(\Delta\eta)^{-(4K-4)}$ (Proposition~\ref{prop:tagged} and Theorem~\ref{thm:discvalue}). Beyond the $K$-cohort case, we prove matching lower and upper bounds of exact logarithmic order $\Theta(1/\log N)$ for general completely monotone memory (Theorem~\ref{thm:general}) and price the resulting disclosure savings against audit-trail budgets. Several ingredients are known in other settings, and we claim only their closed-loop, passive and financial forms: the $\delta^{-4}$ law for telling one component from two closely spaced ones is classical in resolution analysis \citep{shahram2005}; logarithmic rates are known for mixtures of exponentials \citep{rebafka2015} and for general mixing distributions, by the same moment-matching and polynomial-approximation technique \citep{wu2020}; and labelled samples are known to be far more informative than unlabelled ones in mixture problems \citep{castelli1996}.

\paragraph{Estimating heterogeneous-agent models from prices.} A large literature models markets as populations of trader types that switch between forecasting rules \citep{brock1998} and estimates such models from price data. Estimates of the parameters that separate the types are often imprecise: \citet{amilon2008} reports that the estimated model reproduces properties of returns but that its fit is generally poor, \citet{franke2012} compare switching mechanisms by the method of simulated moments, and \citet{kukacka2017} find the switching coefficient statistically insignificant by simulated maximum likelihood while the belief parameters are significant. These papers take the identification of the types from prices as given. Our results suggest one reason the type-separating parameters are hard to pin down: around a steady state, prices determine only the merged measure, and separating types whose memories differ by $\Delta\eta$ takes a number of samples of order $(\Delta\eta)^{-(4K-2)}$ for $K$ types (Theorems~\ref{thm:merged} and~\ref{thm:kcohort}). The models these papers estimate are nonlinear, so this is a statement about their linearization, not a claim that their estimates are wrong; it does give a quantity, the separation of the types' memories, that bounds the precision one can expect before estimating.

\paragraph{Closed-loop system identification.} Engineering has long studied identification under feedback \citep{ljung1999,forssell1999}. There the object is the plant, the difficulty is that feedback correlates the input with the noise, and the remedy is external excitation that makes the experiment informative. Our problem is the reverse in two ways. The object is the feedback law (the traders), not the plant (the price), and the memory rates enter as zeros of the observed transfer function. And external excitation does not remove the difficulty: with a probe the merged measure is identified exactly (Theorem~\ref{thm:merged}), so the barrier is not qualitative identifiability but the conditioning of the map from memories to responses, which we measure in samples, quantity and capital. In finance, \citet{leach1993} show that a market maker can gain by setting prices to learn about demand, so learning through trading has an equilibrium counterpart; our probing observer is a reduced-form version of that idea without the equilibrium. In econometric terms, as $\Delta\eta\to0$ the cohort parameters pass from strong to weak identification \citep{andrews2012}, and our rates state that passage in samples.

\paragraph{Inferring trader types from market data.} Market microstructure has a tradition of inferring hidden trader populations from public data. The probability of informed trading \citep{easley1996} fits a mixture model to daily counts of buys and sells to estimate the share of informed orders, and what it captures has been debated \citep{duarte2009}. With regulatory account data, \citet{kirilenko2017} classify traders by their holding and trading patterns. Our results speak to the case where types differ in how long they remember prices: from prices alone the separation is costly, and account or class records of the kind \citet{kirilenko2017} use remove most of the cost (Proposition~\ref{prop:tagged}).

\begin{table}[t]\centering\small
\begin{threeparttable}
\caption{Comparison with finite-mixture estimation and open-loop identification}\label{tab:novelty}
\setlength{\tabcolsep}{4pt}
\begin{tabular}{*{4}{>{\raggedright\arraybackslash}p{0.225\textwidth}}}
\toprule
Feature & Finite-mixture estimation \citep{heinrich2018,wu2020} & Open-loop exponential fitting \citep{roy1989,batenkov2021} & Closed-loop market identification (this paper)\\
\midrule
Where rates live & Component locations / moments of density & Poles of impulse response & Transmission zeros of closed-loop $H(z)$ (Thm~\ref{thm:merged})\\
Active vs.\ passive data & i.i.d.\ draws only & Designed input vs.\ output spectrum & At least the same $(\Delta\eta)^{-(4K-2)}$ order, uniformly, via the minimum-phase Whittle gap (Thm~\ref{prop:passive})\\
Economic unit of cost & Sample size $N$ & Signal-to-noise ratio & Quantity $\Theta(\delta^{-2})$, trading loss and capital $\Theta(\delta^{-4})$ via $\tfrac12v^\top T_hv$ (Thms~\ref{thm:quantity}--\ref{thm:rh})\\
Effect of class tagging & Observed component labels ($\Theta(1)$) & Separate sensor per mode & Breaks closed loop cohort-by-cohort; saves $\Theta((\Delta\eta)^{-(4K-2)})$ at fixed accuracy, $\Theta((\Delta\eta)^{-(4K-4)})$ for rates (Thm~\ref{thm:discvalue})\\
\bottomrule
\end{tabular}
\end{threeparttable}
\end{table}

\begin{table}[!htbp]\centering\footnotesize
\begin{threeparttable}
\caption{Closest results and how this paper differs}\label{tab:closest}
\setlength{\tabcolsep}{3pt}
\begin{tabular}{*{5}{>{\raggedright\arraybackslash}p{0.172\textwidth}}}
\toprule
Existing result & Observations & Target & Conclusion & Difference from this paper\\
\midrule
Mixture estimation \citep{heinrich2018,wu2020} & i.i.d.\ draws from a $K$-component mixture & Mixing measure & Rates that slow sharply as components merge; $n^{-1/(4K-2)}$ for Gaussian location mixtures & Same exponent; here the measure is observed through a feedback loop, as zeros of $H$, actively or passively, and the cost is also priced in quantity and capital\\
Super-resolution and Prony \citep{moitra2015,batenkov2021,katz2023} & Noisy Fourier or exponential samples & Point-source locations & Stability governed by Vandermonde conditioning and node separation; Prony attains the min-max bound for clustered exponents & Real nodes extracted from a closed loop; our matching upper bound (Theorem~\ref{prop:esprit}) adapts this known result\\
Closed-loop system identification \citep{ljung1999,forssell1999} & Input--output data under feedback & The plant & Consistency requires informative (externally excited) experiments & Target is the feedback law (the traders); excitation gives exact identification of $M$, and the barrier is conditioning, measured in samples and capital\\
Heterogeneous-agent estimation \citep{brock1998,franke2012,kukacka2017} & Prices and returns & Type and switching parameters & Type-separating parameters often imprecise & Bounds that precision by the types' memory separation (in the linearization)\\
Probability of informed trading \citep{easley1996} & Daily buy and sell counts & Share of informed orders & Mixture estimate of trader population & Types that differ in memory; price-only versus class-tagged data\\
Data informativity \citep{vanwaarde2020} & Finite data, often noiseless & Whether a property is determined & Qualitative conditions & Quantitative minimax costs under noise\\
\bottomrule
\end{tabular}
\end{threeparttable}
\end{table}

\paragraph{Statistical identification.} The identification results use two-point testing arguments \citep{lecam1986,tsybakov2009}. The moment-matching construction of Section~\ref{sec:mixtures} is the dynamical analogue of the slow rates in finite-mixture estimation \citep{manole2022} and of super-resolution of near-colliding sources \citep{liu2023,batenkov2021}; the general-memory result combines best-approximation duality \citep{bernstein1914,rivlin1981} with the moment comparisons used by \citet{kong2017} and \citet{wu2020}, and its logarithmic rate parallels supersmooth deconvolution \citep{fan1991}; recovery of exponentials connects to ESPRIT and matrix-pencil methods \citep{roy1989,hua1990}. Minimum-phase spectral factorization and the Whittle likelihood are classical \citep{whittle1953}. What is new here is the mapping from these statistical rates into executed quantity and money through a financially specified account, the passive equivalence, and the value of class records; that the relevant rates are zeros, not poles, of the observable transfer function matters for which estimator is valid, not for the rate.

\paragraph{Non-normal dynamics and supervision.} Transient growth of stable non-normal systems is classical \citep{trefethen2005,kreiss1962}. We use it to show that supervision through eigenvalues and second moments of prices is blind to transient risk (Theorem~\ref{thm:blind}), which motivates the identification-free control of the companion paper. On the regulatory side, pre-trade risk controls for broker-dealers with market access are governed by SEC Rule 15c3-5 \citep{sec2010}, and participant-level audit trails are the natural disclosure regime for Proposition~\ref{prop:tagged}. A large literature studies how pre- and post-trade transparency affects liquidity and trading costs \citep{pagano1996,boehmer2005,bessembinder2006}; we ask a different question, what disclosure lets a supervisor learn about the algorithms that trade, and give the answer in samples. Crypto venues, whose market structure is studied by \citet{makarov2020}, are a natural laboratory for this question because they publish class-level positioning every five minutes, while the CFTC publishes it weekly for regulated futures.

Partial fractions, Gaussian testing, convolution inequalities and positive-system gains are established tools. Our contribution is to connect them into a quantitative chain from memory resolution, through probing resources, to money, with matching orders, explicit constants and a supervisory interpretation, and then to measure the relevant quantities in data.

\begin{table}[t]\centering\small
\begin{threeparttable}
\caption{Map of results}\label{tab:map}
\begin{tabular}{p{0.34\textwidth}p{0.36\textwidth}p{0.18\textwidth}}
\toprule
Question & Answer & Where\\
\midrule
What do prices identify? & The merged rate--weight measure only & Thm~\ref{thm:merged}\\
Quantity to separate a known pair & $\Theta(\delta^{-2})$, exact error formula & Thm~\ref{thm:quantity}\\
Capital to separate a known pair & $\Theta(\delta^{-4})$, counterparty or own funds & Thm~\ref{thm:cap}, Lem.~\ref{thm:capital}\\
Constant at the calibration & $[61.41,64.48]$ par units; $\approx10^5$ USDT & Thms~\ref{thm:spectral}--\ref{thm:rh}, Table~\ref{tab:capital}\\
Does adaptivity help? & At most by a factor $\approx0.49$ in expected loss & Thm~\ref{thm:adaptive}\\
Unknown mixtures & $(\Delta\eta)^{-(4K-2)}$, tight on rich classes for every $K$ & Thms~\ref{thm:kcohort}, \ref{prop:esprit}\\
General memory & Accuracy $\asymp1/\log N$, lower and upper bounds & Lem.~\ref{lem:altpair}, Thm~\ref{thm:general}\\
Passive data & Identifies $M$; at least the same polynomial cost for every $K$ & Thms~\ref{thm:minphase}, \ref{prop:passive}, Prop.~\ref{prop:passive0}\\
Nonlinear and adaptive rules? & Settled online learners obey $\delta^{-4}$; fast adaptation or thresholds break it & Table~\ref{tab:learners}\\
What removes the barrier? & Tagged flow, audited equity & Props~\ref{prop:disclosure}, \ref{prop:tagged}\\
Does attribution change a decision? & Yes in a stylised case: a market-wide limit needs only $S$; whether a targeted halt suffices needs each $q_i$ & Prop.~\ref{prop:target}\\
What is tagging worth? & Factor $(\Delta\eta)^{-(4K-2)}$ (fixed accuracy) or $(\Delta\eta)^{-(4K-4)}$ (rates); $\approx333\times$ for BTCUSDT (illustrative) & Thm~\ref{thm:discvalue}\\
Hypothetical probing loss? & At least \$15M--\$356M a year per market to re-learn two unknown groups weekly at $\Delta\eta=0.10$; illustrative, not a social cost & Table~\ref{tab:cat_dollars}\\
Do class data restrict class persistence? & Yes: CFTC 8 of 8; Binance 30 of 30 and 29 of 29 new markets (registered); whether it is memory is open & Tables~\ref{tab:cot}, \ref{tab:disc}\\
Do the rates survive real price noise? & Yes where it was used directly; large samples use a normal approximation & Table~\ref{tab:injection}\\
Can eigenvalues certify safety? & No & Thm~\ref{thm:blind}\\
What do feedback indices measure? & Mostly filtering and timing; small real feedback & Tables~\ref{tab:panel}, \ref{tab:feedback}\\
Does feedback track activity? & No: fees move volume 2.6--4.9$\times$, not $\hat Q$ & Table~\ref{tab:fee2}\\
\bottomrule
\end{tabular}
\end{threeparttable}
\end{table}

\subsection{Outline}
Section~\ref{sec:model} sets out the model. Section~\ref{sec:ident} gives the identification boundary. Section~\ref{sec:resources} links memory resolution to quantity and money for a known pair. Section~\ref{sec:account} computes account-level capital. Section~\ref{sec:mixtures} treats unknown mixtures and passive data. Section~\ref{sec:learners} checks the results against stylized nonlinear and adaptive trading rules. Section~\ref{sec:disclosure} studies disclosure and compares its cost in dollars with the U.S. Consolidated Audit Trail. Section~\ref{sec:blind} gives the spectral blind spot. Section~\ref{sec:empirics} reports the empirical evidence. Section~\ref{sec:conclusion} concludes. Proofs are in Appendix~\ref{app:proofs}; Appendix~\ref{app:ia} collects additional results, and Appendix~\ref{app:repro} maps every number to its source.

\section{Model and motivating examples}\label{sec:model}

\paragraph{The model in words.} The market has one price and $K$ groups of traders. Each period three things happen. First, the price moves: it drifts back toward its fundamental value, and it is pushed by the groups' trades, by any test trade and by news. Second, each group updates its view by mixing the newest price into its running average, with weight $\eta_i$ on the newest price and $1-\eta_i$ on its old view. Third, each group trades in proportion to its view, with intensity $q_i$. A group with a large $\eta_i$ has a short memory and reacts quickly, like a short-horizon trend follower; a group with a small $\eta_i$ has a long memory and reacts slowly. The equations below write this down.

Time is discrete. Let $p_t$ be the price, $p^*$ the fundamental value and $x_t=p_t-p^*$ the pricing error. A reduced-form dealer sector absorbs net order flow with price impact $\kappa>0$, an imposed primitive rather than an equilibrium quote, and arbitrage pulls the error toward zero with persistence $a\in[0,1)$:
\begin{equation}\label{eq:price}
x_t = a x_{t-1} + \kappa \sum_{i=1}^K q_i y_{i,t-1} + \kappa v_t + w_t,
\end{equation}
where $v_t$ is an executed probe quantity and $w_t$ a price innovation. Cohort $i$ holds a state $y_{i,t}$ with exponential memory and trades $q_i y_{i,t}$:
\begin{equation}\label{eq:cohort}
y_{i,t} = (1-\eta_i)\,y_{i,t-1} + \eta_i x_t,\qquad \eta_i\in(0,1).
\end{equation}
The ordering is \emph{price first}: the price responds to the probe, then cohorts update. Write $r_i=1-\eta_i$ for the memory rate and $w_i=q_i\eta_i$ for the merged weight. The feedback kernel and its transform are
\begin{equation}\label{eq:kernel}
k_m=\sum_i w_i r_i^{\,m},\qquad G(z)=\sum_i \frac{w_i z}{z-r_i},
\end{equation}
and the transfer function from probes to prices, $H(z)=\sum_{j\ge0}h_jz^{-j}$, is
\begin{equation}\label{eq:H}
H(z)=\frac{\kappa z}{z-a-\kappa G(z)}=\frac{\kappa zD(z)}{(z-a)D(z)-\kappa N(z)},\qquad G=\frac ND,\quad D(z)=\prod_{i}(z-r_i),
\end{equation}
with impulse response $h_0=H(\infty)=\kappa,h_1,h_2,\ldots$; in particular $H(1)=\kappa/(1-S)=\sum_jh_j$. Stacking $z_t=(x_t,y_{1,t},\ldots,y_{K,t})^\top$ gives $z_t=Jz_{t-1}+bv_t$, where, because cohorts update on the same period's price, the probe enters every state on impact through
\begin{equation}\label{eq:b}
b=\kappa\,(1,\eta_1,\ldots,\eta_K)^\top ,
\end{equation}
and
\begin{equation}\label{eq:J}
J=\begin{pmatrix} a & \kappa q_1 & \cdots & \kappa q_K\\ \eta_1 a & 1-\eta_1+\eta_1\kappa q_1 & \cdots & \eta_1\kappa q_K\\ \vdots & & \ddots & \vdots\\ \eta_K a & \eta_K\kappa q_1 & \cdots & 1-\eta_K+\eta_K\kappa q_K\end{pmatrix},\qquad S\equiv a+\kappa\sum_i q_i .
\end{equation}
In the \emph{positive economy} ($q_i\ge0$), $J$ is entrywise nonnegative, its price row sums to $S$, and the economy is Schur stable iff $S<1$.

\begin{lemma}[Two motivating examples: filtering and learning]\label{lem:micro}
\begin{enumerate}[label=(\roman*)]
\item \emph{Kalman memory.} A desk that tracks a latent random-walk trend $v_{t+1}=v_t+\zeta_{t+1}$, $\zeta\sim N(0,\sigma_\zeta^2)$, from signals $r_t=v_t+\nu_t$, $\nu\sim N(0,\sigma_\nu^2)$, using the steady-state Kalman filter, updates as $\hat v_{t+1}=(1-\eta)\hat v_t+\eta r_t$ with
\begin{equation}\label{eq:kgain}
\eta=\frac{-\sigma_\zeta^2+\sqrt{\sigma_\zeta^4+4\sigma_\zeta^2\sigma_\nu^2}}{2\sigma_\nu^2}\in(0,1).
\end{equation}
\item \emph{Gradient memory.} Online gradient descent with $L_2$ regularization $\gamma\in(0,1)$, step size $\alpha>0$ and loss $\tfrac12H(x_t-\theta_t)^2$ gives $\theta_{t+1}=(1-\eta)\theta_t+\eta q x_t$ with $\eta=\gamma+\alpha H\in(0,1)$ and $q=\alpha H/\eta$. A tabular value update with learning rate $\eta$ at a fixed state has the same form; general reinforcement-learning policies need not.
\end{enumerate}
\end{lemma}

\begin{remark}[Why exponential memory]
Lemma~\ref{lem:micro} shows that two different design philosophies, Bayesian filtering and regularized online learning, produce the same observable reduced-form belief state $y_{i,t}$. When desk $i$ chooses order flow $u_{i,t}$ subject to $I_{i,t}=(1-\delta_{\rm inv})I_{i,t-1}+u_{i,t}$ and an inventory cap $|I_{i,t}|\le\bar I_i$ to maximize expected continuation return $\beta_i y_{i,t}u_{i,t}$ net of quadratic execution cost $\tfrac12\lambda_{{\rm tc},i}u_{i,t}^2$ and quadratic inventory holding cost $\tfrac12\gamma_{{\rm inv},i}I_{i,t}^2$, its interior first-order condition gives $u_{i,t}^*=q_iy_{i,t}-\psi_{{\rm inv},i}I_{i,t-1}$ with market-response weight $q_i=\beta_i/(\lambda_{{\rm tc},i}+\gamma_{{\rm inv},i})$ and inventory mean-reversion $\psi_{{\rm inv},i}=\gamma_{{\rm inv},i}(1-\delta_{\rm inv})/(\lambda_{{\rm tc},i}+\gamma_{{\rm inv},i})$, which shrinks toward zero as trading costs or inventory penalties tighten and destabilizes the market when the closed-loop spectral radius reaches one (\texttt{microfoundation\_and\_baselines.py}). This filtering-and-control recursion specifies an individual agent's decision rule, not a complete economic equilibrium (dealer--trader equilibrium is analysed in the companion paper).

Two further simplifications matter for reading the results: the observer is taken to know the fundamental value $p^*$, so the pricing error $x_t$ is observed, and the dealer sector is a fixed linear impact $\kappa$ rather than an equilibrium quote. The paragraph before Section~\ref{sec:cot} states what the data tests lose when neither holds. Assuming that the observer knows $p^*$ gives it more information than any real observer has, so the lower bounds hold a fortiori when $p^*$ is unknown; the upper bounds use it.

The model applies these rules to the pricing error $x_t$: a desk that cannot separate fundamental news from transitory deviations treats recent deviations as a signal of a persistent trend, which is the exponential extrapolation of \citet{barberis2015,barberis2018}. With $q_i>0$ such a desk buys after rises, as trend followers do, and nothing in the model says this is profitable; the lemma shows only that the update has the exponential form. In these two illustrative families of algorithms the memory rate is a known function of a tuning parameter, the desk's signal-to-noise ratio or its regularization. That is the kind of private information a supervisor would like to recover, although real trading algorithms need not belong to either family. Section~\ref{sec:ident} shows why prices reveal it only in merged form.
\end{remark}

\begin{lemma}[Positive price-response geometry]\label{lem:cm}
In the positive, stable economy the impulse response is a nonnegative mixture of decaying exponentials,
\[
h_j=\sum_\ell \omega_\ell\lambda_\ell^j,\qquad \omega_\ell\ge0,\quad \sum_\ell\omega_\ell=\kappa,\quad 0\le\lambda_\ell<1 .
\]
It is completely monotone and nonincreasing; its peak is the direct impact $\kappa$, even though $J$ is non-normal. In price/memory coordinates the price-only gain is $\sup_t\|e_x^\top J^t\|_2=1$.
\end{lemma}

As an illustration of the lemma (not part of its proof), $(-1)^n\Delta^nh_j\ge0$ for $n\le4$ held in all 892 stable random positive economies tested. Lemma~\ref{lem:cm} matters twice: it makes probing costs positive (Section~\ref{sec:resources}), and it shows that in this economy transient amplification of the price response cannot arise from cohort feedback alone. The companion paper locates amplification in order-book depth instead.

\subsection{Discussion of the assumptions}\label{sec:assumptions}

\paragraph{Linearity.} The model is linear in prices, flows and memory states. Linearity is what makes the merged-measure equivalence exact and the testing problems Gaussian. Nonlinear strategies (for example threshold rules or position limits) can break the equivalence class, but they do so in regions of the state space that a small probe does not visit; for small probes the linearization around the steady state is the relevant object.

\paragraph{Price-first timing.} We order the price update before the cohort update. Under the alternative lagged timing the transition matrix becomes $\begin{psmallmatrix}a&1\\ \eta Q&1-\eta\end{psmallmatrix}$ and its eigenvalues differ, so stability statements must be computed under the timing actually used. The identification results depend on the transfer function and survive either convention.

\paragraph{Exponential memory.} Lemma~\ref{lem:micro} justifies exponential memory for filtering and regularized learning. Neural and reinforcement-learning policies are not exponential filters. Rules whose orders are smooth functions of an exponential-memory state are exponential filters to first order, so the results hold for small probes; Section~\ref{sec:learners} checks this in stylized simulations and finds where it fails, and how closely production algorithms follow such rules remains an empirical question we do not settle. Mixtures of exponentials can approximate any completely monotone kernel, so a power-law memory corresponds to a continuum of cohorts, for which Theorem~\ref{thm:general} shows that accuracy improves only logarithmically in the data.

\paragraph{Algorithms that react to being probed.} The probing results assume that the cohorts keep their memories while they are probed. An algorithm that detects a probe could pause, shorten its memory or trade against the probe. Four observations limit what this does to our conclusions. First, a reaction that all cohorts share does not help separate them, and an algorithm that mimics another group's memory rate only brings the hypotheses closer, so neither lowers the costs we derive. A reaction that differs across cohorts is a new signature, but one the observer must specify and estimate, so it shifts the identification problem rather than solving it. Second, the passive route needs no probe: Theorem~\ref{prop:passive} shows that natural price variation faces a lower bound of the same polynomial order, so an observer who avoids detection by not trading does not escape the barrier; it pays in time instead of money. We do not show that passive data attain that order. Third, detection is itself a testing problem against the same noise: the generalized-eigenvector probe of Section~\ref{sec:account} trades at most 1.25 par units in any period, and a probe sized within normal flow variation is hard to tell from noise, so the threat bites mainly for large single pulses, which our ledger shows are not optimal anyway. Fourth, what we cannot rule out is an algorithm designed to defeat surveillance, which turns identification into a game between observer and observed; our results are the benchmark against which such a game would be measured, and we leave it to future work.

\paragraph{Known primitives.} The resource theorems take $\kappa$, $a$ and the noise law as known. Unknown primitives can only make separation harder, so the lower bounds remain valid. The upper bounds as stated assume the primitives known; estimating them adds a cost we do not analyse, which we expect to be of lower order because they enter the response at first order. In an equilibrium in which dealers set the impact $\kappa$ from the order flow they see, as in \citet{kyle1985}, a prober's unexpected trades would also move $\kappa$; we do not model this, and expect it to raise the cost of probing rather than lower it.

\paragraph{Observation contracts.} Each result names what the observer sees. A price-only observer, an observer with executed quantities, one with noisy statements, one with audited equity and one with tagged flow face different experiments. The paper's central message is that these differences, not the observer's statistical skill, determine what can be learned.
\subsection{The observer, its options and its costs}\label{sec:observer}
The results are stated from the side of one observer outside the trader groups: a supervisor, an exchange's surveillance unit or a large desk. Its \emph{target} is attribution, recovering the merged weight and memory rate of each group to a stated accuracy, or deciding between two candidate populations at a stated error. Before it observes, it knows the price persistence $a$, the impact $\kappa$ and the law of the observation noise, and, for the $K$-group results, the number of groups; it does not know their rates or weights. These are the conditions under which the lower bounds are statements about the data rather than about the observer's ignorance; unknown primitives can only enlarge the set of candidate populations.

The observer has three options, and each result prices one of them.
\begin{enumerate}[label=(\alph*),itemsep=1pt]
\item \emph{Wait.} Record stationary prices driven by natural innovations, without trading. The cost is time: the number of observations needed (Theorems~\ref{thm:minphase} and~\ref{prop:passive}).
\item \emph{Probe.} Execute a sequence $v_t$ of trades against the groups and record the price response. Probes execute at the average price of a linear book after the uncontrolled state update; the prober bears the execution loss $\tfrac12v^\top T_hv$ plus fees (Lemma~\ref{lem:toeplitz}) and must fund margin throughout. The counted resource is the number of independent experiments, each starting from the reset state, or the executed quantity and capital of a single experiment (Theorems~\ref{thm:quantity}--\ref{thm:spectral}).
\item \emph{Acquire disclosure.} Receive each class's flow or position, observed with noise, at an operating cost per supervisory cycle and a compliance cost per market (Proposition~\ref{prop:tagged}, Theorem~\ref{thm:discvalue}, Proposition~\ref{prop:choice}).
\end{enumerate}
Within each period the order is: the probe (if any) executes and the price updates, the groups update their states on the new price, and the observer records the price with noise. The groups do not change their rules in response to being observed; Section~\ref{sec:assumptions} discusses what happens if they do. Probing experiments start from the reset state. A probe sent into a running market adds the free response of the current state, which by linearity superposes with the forced response; if the state is unknown it acts as a nuisance term, which can only make separation harder, so the lower bounds still hold, while the upper bounds as stated need the reset state or a known one. Independent resets are a benchmark experiment that no outside participant can run in a live market; we use them because they make the information content of the data exact, and the economic comparisons say whenever they rely on them. Regulators do not, in practice, trade to learn; the probing option describes an outside participant, or prices the information a supervisor would otherwise have to buy by trading.

Three money quantities appear and are kept apart: the \emph{capital required}, the equity an account must hold so that its margin constraint is never breached while it probes (Theorem~\ref{thm:cap} and Lemma~\ref{thm:capital}); the \emph{round-trip trading loss}, the execution loss of a probe that ends flat, which is what the account actually spends (Theorems~\ref{thm:spectral}, \ref{thm:adaptive} and~\ref{thm:rh}); and the \emph{annualised loss} of renewing that probe weekly, used for the comparison with audit-trail budgets (Section~\ref{sec:disc_dollars}). The lower bounds hold for the stated probe classes, deterministic probes and, in Theorem~\ref{thm:adaptive}, non-anticipating sequential probes; they are not lower bounds on information obtained through other channels, which is the point of Section~\ref{sec:disclosure}. Two rates recur and measure different things: $\delta^{-4}$ is the capital (and the average-power horizon) for a \emph{known} symmetric pair whose rates differ by $2\delta$, while $(\Delta\eta)^{-(4K-2)}$, which is $(\Delta\eta)^{-6}$ for $K=2$, is the number of experiments for $K$ \emph{unknown} groups with spacing $\Delta\eta$.

